\documentclass{amsart}
% For dealing with references we use the comment environment
\usepackage{verbatim}
\newenvironment{reference}{\comment}{\endcomment}
%\newenvironment{reference}{}{}
\newenvironment{slogan}{\comment}{\endcomment}
\newenvironment{history}{\comment}{\endcomment}

% For commutative diagrams we use Xy-pic
\usepackage[all]{xy}

% We use 2cell for 2-commutative diagrams.
\xyoption{2cell}
\UseAllTwocells

% We use multicol for the list of chapters between chapters
\usepackage{multicol}

% This is generally recommended for better output
\usepackage{lmodern}
\usepackage[T1]{fontenc}

% For cross-file-references

% Package for hypertext links:
\usepackage{hyperref}

% For any local file, say "hello.tex" you want to link to please

% Theorem environments.
%
\theoremstyle{plain}
\newtheorem{theorem}[subsection]{Theorem}
\newtheorem{proposition}[subsection]{Proposition}
\newtheorem{lemma}[subsection]{Lemma}

\theoremstyle{definition}
\newtheorem{definition}[subsection]{Definition}
\newtheorem{example}[subsection]{Example}
\newtheorem{exercise}[subsection]{Exercise}
\newtheorem{situation}[subsection]{Situation}

\theoremstyle{remark}
\newtheorem{remark}[subsection]{Remark}
\newtheorem{remarks}[subsection]{Remarks}

\numberwithin{equation}{subsection}

% Macros
%
\def\lim{\mathop{\mathrm{lim}}\nolimits}
\def\colim{\mathop{\mathrm{colim}}\nolimits}
\def\Spec{\mathop{\mathrm{Spec}}}
\def\Hom{\mathop{\mathrm{Hom}}\nolimits}
\def\Ext{\mathop{\mathrm{Ext}}\nolimits}
\def\SheafHom{\mathop{\mathcal{H}\!\mathit{om}}\nolimits}
\def\SheafExt{\mathop{\mathcal{E}\!\mathit{xt}}\nolimits}
\def\Sch{\mathit{Sch}}
\def\Mor{\mathop{\mathrm{Mor}}\nolimits}
\def\Ob{\mathop{\mathrm{Ob}}\nolimits}
\def\Sh{\mathop{\mathit{Sh}}\nolimits}
\def\NL{\mathop{N\!L}\nolimits}
\def\CH{\mathop{\mathrm{CH}}\nolimits}
\def\proetale{{pro\text{-}\acute{e}tale}}
\def\etale{{\acute{e}tale}}
\def\QCoh{\mathit{QCoh}}
\def\Ker{\mathop{\mathrm{Ker}}}
\def\Im{\mathop{\mathrm{Im}}}
\def\Coker{\mathop{\mathrm{Coker}}}
\def\Coim{\mathop{\mathrm{Coim}}}

% Boxtimes
%
\DeclareMathSymbol{\boxtimes}{\mathbin}{AMSa}{"02}

%
% Macros for moduli stacks/spaces
%
\def\QCohstack{\mathcal{QC}\!\mathit{oh}}
\def\Cohstack{\mathcal{C}\!\mathit{oh}}
\def\Spacesstack{\mathcal{S}\!\mathit{paces}}
\def\Quotfunctor{\mathrm{Quot}}
\def\Hilbfunctor{\mathrm{Hilb}}
\def\Curvesstack{\mathcal{C}\!\mathit{urves}}
\def\Polarizedstack{\mathcal{P}\!\mathit{olarized}}
\def\Complexesstack{\mathcal{C}\!\mathit{omplexes}}
% \Pic is the operator that assigns to X its picard group, usage \Pic(X)
% \Picardstack_{X/B} denotes the Picard stack of X over B
% \Picardfunctor_{X/B} denotes the Picard functor of X over B
\def\Pic{\mathop{\mathrm{Pic}}\nolimits}
\def\Picardstack{\mathcal{P}\!\mathit{ic}}
\def\Picardfunctor{\mathrm{Pic}}
\def\Deformationcategory{\mathcal{D}\!\mathit{ef}}

\setlength{\emergencystretch}{3em}
\begin{document}
\title[Stacks Project, Tag 08PX]{393 --- Finite polynomial simplicial resolutions of finite-type algebras}
\maketitle
\noindent Copyright (C) 2005--2025 Johan de Jong. Stacks Project authors. Source: \href{https://stacks.math.columbia.edu/tag/08PX}{Tag 08PX}.

\noindent Editorial difficulty note (not author text). Difficulty H2: The proof inductively replaces a coskeletal resolution in its next degree by a finite polynomial algebra. It must preserve the trivial Kan fibration through the replacement and prove that the inverse-limit resolution stabilizes degreewise. This is an explicit higher resolution construction, rather than an application of abstract resolution existence..

\noindent Editorial provenance note (not author text). History: excerpt from revision \texttt{a0444\allowbreak 6e57e\allowbreak c1fbc\allowbreak 252a8\allowbreak 71afc\allowbreak ec775\allowbreak 2fb28\allowbreak 07b14}, cotangent.tex, lines 554--635. Reference presentation was adapted; original-author context and core support blocks were retained or added. The mathematical wording is unchanged. Licensed under GNU FDL 1.2 or later; no invariant sections or cover texts. See ../licenses/COPYING. Context blocks reproduce author definitions and supporting results; they are not additional counted problems.

\begin{definition}
\label{definition-standard-resolution}
Let $A \to B$ be a ring map. The {\it standard resolution of $B$ over $A$}
is the augmentation $\epsilon : P_\bullet \to B$ with terms
$$
P_0 = A[B],\quad P_1 = A[A[B]],\quad \ldots
$$
and maps as constructed in
Simplicial, Example \href{https://stacks.math.columbia.edu/tag/09CB}{example polynomial algebra maps (Tag 09CB)}.
\end{definition}

\section{Truncated simplicial objects and skeleton functors}
\label{simplicial-section-skeleton}

\noindent
Let $\Delta_{\leq n}$ denote the full subcategory of
$\Delta$ with objects $[0], [1], [2], \ldots, [n]$.
Let $\mathcal{C}$ be a category.



\section{Coskeleton functors}
\label{simplicial-section-coskeleton}

\noindent
Let $\mathcal{C}$ be a category.
The {\it coskeleton functor} (if it exists) is a functor
$$
\text{cosk}_n :
\text{Simp}_n(\mathcal{C}) \longrightarrow \text{Simp}(\mathcal{C})
$$
which is right adjoint to the skeleton functor. In a formula
\begin{equation}
\label{simplicial-equation-cosk}
\Mor_{\text{Simp}(\mathcal{C})}(U, \text{cosk}_n V)
=
\Mor_{\text{Simp}_n(\mathcal{C})}(\text{sk}_n U, V)
\end{equation}
Given a $n$-truncated simplicial object $V$ we
say that {\it $\text{cosk}_nV$ exists} if there
exists a $\text{cosk}_nV \in \Ob(\text{Simp}(\mathcal{C}))$
and a morphism $\text{sk}_n \text{cosk}_n V \to V$
such that the displayed formula holds, in other words
if the functor
$U \mapsto \Mor_{\text{Simp}_n(\mathcal{C})}(\text{sk}_n U, V)$
is representable. If it exists it
is unique up to unique isomorphism by the Yoneda lemma.
See Categories, Section \href{https://stacks.math.columbia.edu/tag/001L}{section opposite (Tag 001L)}.



\begin{proposition}
\label{proposition-polynomial}
Let $A$ be a Noetherian ring. Let $A \to B$ be a finite type ring map.
There exists a simplicial $A$-algebra $P_\bullet$ with an augmentation
$\epsilon : P_\bullet \to B$ such that each $P_n$ is a polynomial algebra
of finite type over $A$ and such that $\epsilon$ is a trivial
Kan fibration of simplicial sets.
\end{proposition}

\begin{proof}
Let $\mathcal{A}$ be the category of $A$-algebra maps $C \to B$.
In this proof our simplicial objects and skeleton and coskeleton
functors will be taken in this category.

\medskip\noindent
Choose a polynomial algebra $P_0$ of finite type over $A$ and a surjection
$P_0 \to B$. As a first approximation we take
$P_\bullet = \text{cosk}_0(P_0)$. In other words, $P_\bullet$ is the simplicial
$A$-algebra with terms $P_n = P_0 \times_A \ldots \times_A P_0$.
(In the final paragraph of the proof this simplicial object will
be denoted $P^0_\bullet$.) By
Simplicial, Lemma \href{https://stacks.math.columbia.edu/tag/01AB}{lemma cosk minus one equivalence (Tag 01AB)}
the map $P_\bullet \to B$ is a trivial Kan fibration of simplicial sets.
Also, observe that $P_\bullet = \text{cosk}_0 \text{sk}_0 P_\bullet$.

\medskip\noindent
Suppose for some $n \geq 0$ we have constructed $P_\bullet$
(in the final paragraph of the proof this will be $P^n_\bullet$)
such that
\begin{enumerate}
\item[(a)] $P_\bullet \to B$ is a trivial Kan fibration of simplicial sets,
\item[(b)] $P_k$ is a finitely generated polynomial algebra for
$0 \leq k \leq n$, and
\item[(c)] $P_\bullet = \text{cosk}_n \text{sk}_n P_\bullet$
\end{enumerate}
By Lemma \href{https://stacks.math.columbia.edu/tag/08PW}{lemma polynomial (Tag 08PW)}
we can find a finitely generated polynomial algebra $Q$ over $A$
and a surjection $Q \to P_{n + 1}$. Since $P_n$ is a polynomial algebra
the $A$-algebra maps $s_i : P_n \to P_{n + 1}$ lift to maps
$s'_i : P_n \to Q$. Set $d'_j : Q \to P_n$ equal to the composition of
$Q \to P_{n + 1}$ and $d_j : P_{n + 1} \to P_n$.
We obtain a truncated simplicial object $P'_\bullet$ of $\mathcal{A}$
by setting $P'_k = P_k$ for $k \leq n$ and $P'_{n + 1} = Q$ and morphisms
$d'_i = d_i$ and $s'_i = s_i$ in degrees $k \leq n - 1$ and using the
morphisms $d'_j$ and $s'_i$ in degree $n$. Extend this to a full simplicial
object $P'_\bullet$ of $\mathcal{A}$ using $\text{cosk}_{n + 1}$. By
functoriality of the coskeleton functors there is a morphism
$P'_\bullet \to P_\bullet$ of simplicial objects extending the
given morphism of $(n + 1)$-truncated simplicial objects.
(This morphism will be denoted $P^{n + 1}_\bullet \to P^n_\bullet$
in the final paragraph of the proof.)

\medskip\noindent
Note that conditions (b) and (c) are satisfied for $P'_\bullet$ with $n$
replaced by $n + 1$. We claim the map $P'_\bullet \to P_\bullet$ satisfies
assumptions (1), (2), (3), and (4) of
Simplicial, Lemmas \href{https://stacks.math.columbia.edu/tag/01A6}{lemma section (Tag 01A6)}
with $n + 1$ instead of $n$. Conditions (1) and (2) hold by construction.
By Simplicial, Lemma \href{https://stacks.math.columbia.edu/tag/08NJ}{lemma cosk above object (Tag 08NJ)}
we see that we have
$P_\bullet = \text{cosk}_{n + 1}\text{sk}_{n + 1}P_\bullet$
and
$P'_\bullet = \text{cosk}_{n + 1}\text{sk}_{n + 1}P'_\bullet$
not only in $\mathcal{A}$ but also in the category of $A$-algebras,
whence in the category of sets (as the forgetful functor from $A$-algebras
to sets commutes with all limits). This proves (3) and (4). Thus the lemma
applies and $P'_\bullet \to P_\bullet$ is a trivial Kan fibration. By
Simplicial, Lemma \href{https://stacks.math.columbia.edu/tag/08NP}{lemma trivial kan composition (Tag 08NP)}
we conclude that $P'_\bullet \to B$ is a trivial Kan fibration and (a)
holds as well.

\medskip\noindent
To finish the proof we take the inverse limit $P_\bullet = \lim P^n_\bullet$
of the sequence of simplicial algebras
$$
\ldots \to P^2_\bullet \to P^1_\bullet \to P^0_\bullet
$$
constructed above. The map $P_\bullet \to B$ is a trivial Kan fibration by
Simplicial, Lemma \href{https://stacks.math.columbia.edu/tag/08NQ}{lemma limit trivial kan (Tag 08NQ)}.
However, the construction above stabilizes in each degree
to a fixed finitely generated polynomial algebra as desired.
\end{proof}
\end{document}
